% Motivation and learning outcomes.
% The outcomes are the outcomes of Day 8 in the syllabus.
\begin{frame}{Why this day matters}
  \begin{itemize}
    \item Days 5 and 7 gave one label for a complete image. A detector answers
          two questions for each object: which class it has, and where it is.
          Its output is a list of boxes, each with a class and a score.
    \item A single-stage detector gives thousands of candidate boxes for one
          image. The post-processing removes the duplicates. This step runs on
          the processor after the model, and it is a part of the latency.
    \item One accuracy value does not describe a detector. IoU, precision,
          recall, and mAP say how good the boxes are. The input resolution
          changes the frame rate and the smallest object that the model finds.
    \item In the lab you run an SSD model and a YOLO model on the Raspberry
          Pi, train a YOLO model on a small custom dataset, export it to NCNN
          and to LiteRT in int8, and measure the frame rate with the camera.
  \end{itemize}
\end{frame}

\begin{frame}{Learning outcomes}
  After this day you can:
  \begin{itemize}
    \item Describe the structure of a single-stage detector.
    \item Compute IoU, apply non-maximum suppression, and read an mAP result.
    \item Train a custom YOLO model and export it for the Raspberry Pi.
    \item Tune input resolution and precision for a target frame rate.
  \end{itemize}
\end{frame}
